\documentclass[12pt, letterpaper]{article}

\title{CS302: Testing \color{red}\textbf{[KEY]}\normalcolor}
\author{Amiel L. Iliesi}
\date{2026-04-08}

% packages
\usepackage[hidelinks]{hyperref}
\usepackage[x11names]{xcolor}
\usepackage{minted}

% new commands
\newcommand{\question}[1]{\textbf{#1}\par}
\newcommand{\answer}[1]{\color{Green4}\textbf{#1}\normalcolor}

% new environments
\newenvironment{answerblock}{\color{Green4}}{\normalcolor\vspace{1em}}

% global settings
\setlength{\parindent}{0pt}
\setlength{\parskip}{0.5em}
\setminted{bgcolor=darkgray, style=monokai, fontsize=\footnotesize}

\begin{document}

\maketitle

\tableofcontents

\pagebreak
\section{Testing Data Structures}

\question{In the following sections we will outline a few ways to test data
structures. What do the tests tell us about our data structure, if they all
pass?}

\begin{answerblock}
They tell us that our data structure has all the properties we expect--with
confidence. If we couldn't be sure of our data structure, bugs would emerge
downstream all the time when using our data structures.
\end{answerblock}

\question{What might happen if we forgo testing our data structures? What
benefits does testing give to us programmers \emph{while} we are building the
data structure?}

\begin{answerblock}
If we don't test our data structures, we obviously lose our confidence of data
integrity, and cannot rely on the benefits of our unique data structure, based
on its now-invalid properties. On top of that, as developers of the structure,
we miss catching bugs early that might create compounding bugs that depend on
structural invariants being true. IE if you don't test properties early on,
then random bugs will pop up later that seem mystifying, and become much harder
to debug.
\end{answerblock}

\pagebreak
\section{Blackbox}

\question{Describe blackbox testing. What does it test, what is known while
testing?}

\begin{answerblock}
Blackbox testing is one of our two methods we use to test our data structures.
It is called blackbox, because we pretend (or enforce) our view of our data
structure to be black--or, opaque. We can test our structure's function return
values in order to assess whether or not the structure has been operating
correctly.

We only know the things a user would know about the function, public members
and such.
\end{answerblock}

\question{If you wanted to blackbox test a balanced binary tree, what are some
tests you could implement?}

\begin{answerblock}
You could test:

\begin{itemize}
    \item Whether or not all inserted items made it in.
    \item Whether or not deleting reduces the number of elements correctly.
\end{itemize}

Or any other tests that validate data integrity.
\end{answerblock}

\pagebreak
\section{Glassbox}

\question{Describe glassbox testing. What does it test, what is known while
testing?}

\begin{answerblock}
Glassbox testing is a type of testing that assumes we can look into our
structure (like it's made of glass). This means that we can verify the
structure itself, instead of just data integrity like with blackbox testing.

With glassbox testing, we know everything about the structure.
\end{answerblock}

\question{If you wanted to glassbox test a balanced binary tree, what are some
tests you could implement?}

\begin{answerblock}
You could test:

\begin{itemize}
    \item Whether or not the tree is balanced (left height \(\approx\) right height).
    \item Whether or not all elements are less or greater than parents
    depending on their direction.
    \item If we have parent pointers, that all parent and children pointers are
    bidirectional.
\end{itemize}

Basically, every invariant/property of our class can have one or more tests to
assure that property with glassbox testing.
\end{answerblock}

\pagebreak
\section{Unit Testing}

\question{What is the point of unit testing? Describe an environment where you
would see this, and the purpose it fulfills.}

\begin{answerblock}
Unit testing is great if you are writing a large project, or if you have a
large project environment with multiple contributors, and you want to ensure
that function domain and ranges are verified. It makes sure that new code is
written on top of predictable input/output to try and reduce bugs and
tech-debt.
\end{answerblock}

\question{Write a basic unit test, and show it's execution. Do this
however you like, C++, Python, pseudocode, just be specific.}

\begin{minted}{python}
def auth_user(user: User, project: Project) -> bool:
    return user in project.editors and user.status is User.Status.ACTIVE

def auth_user_test_1() -> bool:
    user = User(status=User.Status.FROZEN)
    project = Project(editors=[user])
    return auth_user(user, project) == False

def auth_user_test_2() -> bool:
    user = User(status=User.Status.ACTIVE)
    project = Project(editors=[user])
    return auth_user(user, project) == True

# ...etc
# ...

test_results = []
for test_func in tests: # pretend `tests` is all auth tests 
    curr = test_func()
    test_results.append(curr)
    if curr == False:
        print(f'*** {test_func.__name__} fails')

if all(test_results):
    print('*** all tests passed')
\end{minted}

\begin{minipage}{\textwidth}
\question{BONUS: do you know the other big two types of software testing?}

\begin{answerblock}
In addition to unit testing, there's \textbf{integration testing}, and
\textbf{end-to-end}/\textbf{system testing}. Integration testing seeks to test
how components interact with one another, rather then testing individual
components. End-to-end is the next and furthest level out. It tests all
components together, alongside user stories, UX, and other product-wide
metrics.
\end{answerblock}
\end{minipage}

\end{document}